\chapter{Voraussetzungen}

\section{Grundlagen der Mathematik}

\subsection{Notation}

\begin{table}[h!]
    \centering
    \begin{tabularx}{\textwidth}{lX}
        \textbf{Symbol} & \textbf{Bedeutung} \\\toprule
        $x$                        & Skalar (kursiver Kleinbuchstabe) \\
        $\mathbf{x}$               & Vektor (fetter Kleinbuchstabe, Spaltenvektor) \\
        $\mathbf{x}_{i}$             & Vektorelement (i-tes Element) \\
        $\mathbf{x}^\top$          & Transponierter Zeilenvektor \\
        $\mathbf{W}$               & Matrix (fetter Großbuchstabe) \\
        $W_{i,j}$                   & Element in Zeile $i$, Spalte $j$ von $\mathbf{W}$ \\
        $\mathbb{R}$               & Menge der reellen Zahlen \\
        $\mathbb{R}^n$             & $n$-dimensionaler reeller Vektorraum \\
        $\odot$                    & Hadamard-Produkt (elementweise Multiplikation) \\\bottomrule
    \end{tabularx}
    \caption{Im Dokument verwendete mathematische Notation}
    \label{tab:notation}
\end{table}

\subsection{Partielle Ableitung}

Gegeben eine Funktion $f : \mathbb{R}^n \to \mathbb{R}$, die von mehreren Variablen
$x_1, x_2, \ldots, x_n$ abhängt, beschreibt die partielle Ableitung
\begin{equation}
    \frac{\partial f}{\partial x_i}
\end{equation}
die Änderungsrate von $f$ bezüglich $x_i$, während alle anderen Variablen konstant gehalten werden.
Partielle Ableitungen sind die Grundlage der mehrdimensionalen Differentialrechnung und bilden
das mathematische Fundament für Optimierungsverfahren im maschinellen Lernen.

\subsection{Kettenregel}

Die Kettenregel beschreibt, wie die Ableitung einer zusammengesetzten Funktion berechnet wird.
Sind $f$ und $g$ differenzierbare Funktionen mit $h(x) = f(g(x))$, so gilt:
\begin{equation}
    \frac{\mathrm{d}h}{\mathrm{d}x} = \frac{\mathrm{d}f}{\mathrm{d}g} \cdot \frac{\mathrm{d}g}{\mathrm{d}x}
\end{equation}

\subsection{Gradient}

Der Gradient einer skalaren Funktion $f : \mathbb{R}^n \to \mathbb{R}$ ist der Vektor aller
partiellen Ableitungen:
\begin{equation}
    \nabla f(\mathbf{x}) =
    \begin{pmatrix}
        \dfrac{\partial f}{\partial x_1} \\[6pt]
        \vdots \\[2pt]
        \dfrac{\partial f}{\partial x_n}
    \end{pmatrix}
\end{equation}
Er zeigt in die Richtung des steilsten Anstiegs von $f$ am Punkt $\mathbf{x}$.

\subsection{Jacobi-Matrix}

Die Jacobi-Matrix verallgemeinert den Gradienten auf Funktionen mit Vektoren.
Gegeben eine Funktion $\mathbf{f} : \mathbb{R}^n \to \mathbb{R}^m$, die einen Eingangsvektor
$\mathbf{x} \in \mathbb{R}^n$ auf einen Ausgangsvektor $\mathbf{f}(\mathbf{x}) \in \mathbb{R}^m$ abbildet,
ist die Jacobi-Matrix $\mathbf{J} \in \mathbb{R}^{m \times n}$ definiert als:
\begin{equation}
    \mathbf{J} = \frac{\partial \mathbf{f}}{\partial \mathbf{x}} =
    \begin{pmatrix}
        \dfrac{\partial f_1}{\partial x_1} & \cdots & \dfrac{\partial f_1}{\partial x_n} \\[8pt]
        \vdots & \ddots & \vdots \\[4pt]
        \dfrac{\partial f_m}{\partial x_1} & \cdots & \dfrac{\partial f_m}{\partial x_n}
    \end{pmatrix}.
\end{equation}
Das Element $J_{ij} = \frac{\partial f_i}{\partial x_j}$ gibt an, wie stark sich die $i$-te
Ausgabekomponente ändert, wenn die $j$-te Eingabekomponente variiert wird.

Für den Spezialfall $m = 1$ reduziert sich die Jacobi-Matrix zum transponierten Gradienten:
$\mathbf{J} = \nabla f^\top$.

\subsection{Hadamard-Produkt}

Das Hadamard-Produkt, symbolisiert durch $\odot$, bezeichnet die elementweise Multiplikation von zwei Matrizen oder Vektoren gleicher Dimension.
Gegeben zwei Matrizen $\mathbf{A}, \mathbf{B} \in \mathbb{R}^{m \times n}$, ist ihr Hadamard-Produkt $\mathbf{C} = \mathbf{A} \odot \mathbf{B}$ eine Matrix aus $\mathbb{R}^{m \times n}$, deren Elemente wie folgt definiert sind:
\begin{equation}
    C_{i,j} = A_{i,j} \cdot B_{i,j}
\end{equation}
Für Vektoren $\mathbf{a}, \mathbf{b} \in \mathbb{R}^n$ berechnet sich das Hadamard-Produkt äquivalent als $c_i = a_i \cdot b_i$.

% -------------------------------------------------------------------------

\section{Grundlagen des Maschinellen Lernens}

\subsection{Datenvorbereitung}

Bevor ein Modell trainiert werden kann, müssen die Rohdaten in eine geeignete Form gebracht werden.
Dieser Schritt wird als Datenvorbereitung bezeichnet und umfasst mehrere Teilaufgaben.

Zunächst werden fehlende Werte behandelt, etwa durch Entfernen betroffener Einträge oder durch
das Einsetzen von statistischen Kennzahlen wie dem Mittelwert.
Anschließend erfolgt häufig eine Normalisierung oder Standardisierung der Eingabemerkmale,
damit numerisch sehr unterschiedliche Größen den Trainingsprozess nicht verschlechtern.
Bei der Standardisierung wird jedes Merkmal $x$ transformiert zu
\begin{equation}
    x' = \frac{x - \mu}{\sigma},
\end{equation}
wobei $\mu$ der Mittelwert und $\sigma$ die Standardabweichung des jeweiligen Merkmals sind.

Schließlich wird der Datensatz in mindestens zwei disjunkte Teilmengen aufgeteilt:
einen Trainings- und einen Testdatensatz. Optional wird ein dritter Validierungsdatensatz
abgetrennt, der zur Hyperparameterauswahl während des Trainings genutzt wird.

\subsection{Regression}

Regression bezeichnet statistische Methoden zur Modellierung der Beziehung zwischen einer Zielgröße
und einer oder mehreren unabhängigen Merkmalen.
Das einfachste Modell ist die lineare Regression, die eine lineare Abbildung von der
Eingabe $\mathbf{x} \in \mathbb{R}^n$ auf eine vorhergesagte Ausgabe $\hat{y} \in \mathbb{R}$ annimmt:
\begin{equation}
    \hat{y} = \mathbf{w}^\top \mathbf{x} + b,
\end{equation}
wobei $\mathbf{w}$ der Gewichtsvektor und $b$ der Bias-Term sind.

Die Qualität der Vorhersage wird durch eine Verlustfunktion $\mathcal{L}$ gemessen.
Das Ziel des Lernalgorithmus ist es, die Parameter $\mathbf{w}$ und $b$ so zu wählen, dass
$\mathcal{L}$ minimiert wird.

\subsection{Neuronale Netze}

Ein künstliches neuronales Netz besteht aus mehreren hintereinandergeschalteten Schichten
von sog. Neuronen. Jedes Neuron berechnet eine gewichtete Summe seiner Eingaben und wendet
darauf eine nichtlineare Aktivierungsfunktion $\sigma$ an:
\begin{equation}
    a = \sigma\!\left(\mathbf{w}^\top \mathbf{x} + b\right).
\end{equation}

Eine Schicht transformiert einen Eingangsvektor $\mathbf{x}$ zu einem Ausgangsvektor $\mathbf{a}$
mittels einer Gewichtsmatrix $\mathbf{W}$ und eines Biasvektors $\mathbf{b}$:
\begin{equation}
    \mathbf{a} = \sigma\!\left(\mathbf{W}\mathbf{x} + \mathbf{b}\right).
\end{equation}
Durch das Hintereinanderschalten mehrerer solcher Schichten entsteht ein
\textit{tiefes} Netz (\textit{Deep Neural Network}), das in der Lage ist, hochgradig
nichtlineare Zusammenhänge zu erlernen.
Das gesamte Netz realisiert damit eine parametrisierte Funktion
$f_\theta : \mathbb{R}^n \to \mathbb{R}^m$, deren Parameter $\theta = \{\mathbf{W}^{*}, \mathbf{b}^{*}\}$
durch Training bestimmt werden.

\subsection{Verlustfunktion}

Die Verlustfunktion (auch Loss-Funktion oder Kostenfunktion genannt) quantifiziert die Abweichung zwischen der vom Modell vorhergesagten Ausgabe $\mathbf{\hat{y}}$ und dem tatsächlichen Zielwert $\mathbf{y}$. Die Minimierung dieser Funktion $\mathcal{L}$ während des Trainingsprozesses ist das Hauptziel von Optimierungsalgorithmen im maschinellen Lernen \autocite{Goodfellow2016}.

Abhängig von der zugrundeliegenden Problemstellung kommen üblicherweise unterschiedliche Verlustfunktionen $l$ zum Einsatz \autocite{Bishop2006}.

Die allgemeine Lossfunktion für einen Datensatz mit $n$ Beispielen wird durch den Durchschnitt der Einzelfehler über alle Beispiele definiert:
\begin{equation}
    \mathcal{L} = \frac{1}{n} \
    \sum_{i=1}^n l\!\left(\mathbf{\hat{y}}_i, \mathbf{y}_i\right)
\end{equation}

\subsection{Backpropagation}

Der Algorithmus zur Berechnung der Gradienten in neuronalen Netzen wird als
Backpropagation bezeichnet. Er wendet die Kettenregel der Differentialrechnung rekursiv an,
um die partiellen Ableitungen der Verlustfunktion $\mathcal{L}$ nach allen Parametern
effizient zu berechnen.

In der Praxis wird Backpropagation mit einem Gradientenabstiegsverfahren kombiniert.
Beim stochastischen Gradientenabstieg (\textit{Stochastic Gradient Descent}, SGD) werden
die Parameter nach jedem Mini-Batch aktualisiert:
\begin{equation}
    (\mathbf{W}, \mathbf{b}) \leftarrow (\mathbf{W}, \mathbf{b}) - \eta \, \nabla_{(\mathbf{W}, \mathbf{b})} \mathcal{L}
\end{equation}
wobei $\eta > 0$ die Lernrate ist.
// TODO: Weitere Informationen im Abschnitt anfügen
